\documentclass[12pt,a4paper]{amsart}

\usepackage{amsmath,amssymb,amsthm,mathtools,enumitem,booktabs,array,hyperref}
\usepackage[utf8]{inputenc}
\usepackage[T1]{fontenc}
\usepackage[ngerman]{babel}
\usepackage{geometry}
\geometry{margin=2.5cm}

% -------------------------------------------------------
% Theorem-Umgebungen (Englisch)
% -------------------------------------------------------
\newtheorem{theorem}{Theorem}[section]
\newtheorem{lemma}[theorem]{Lemma}
\newtheorem{corollary}[theorem]{Korollar}
\newtheorem{proposition}[theorem]{Proposition}
\theoremstyle{definition}
\newtheorem{definition}[theorem]{Definition}
\newtheorem{example}[theorem]{Beispiel}
\newtheorem{remark}[theorem]{Bemerkung}
\newtheorem{conjecture}[theorem]{Conjecture}

% Deutschsprachige Theorem-Umgebungen
\newtheorem{satz}[theorem]{Satz}
\newtheorem{vermutung}[theorem]{Vermutung}
\newtheorem{korollar}[theorem]{Korollar}
\newtheorem{bemerkung}[theorem]{Bemerkung}
\newtheorem{definition_de}[theorem]{Definition}

% -------------------------------------------------------
% Mathematische Abkürzungen
% -------------------------------------------------------
\newcommand{\N}{\mathbb{N}}
\newcommand{\Z}{\mathbb{Z}}
\newcommand{\Q}{\mathbb{Q}}
\newcommand{\R}{\mathbb{R}}
\newcommand{\C}{\mathbb{C}}
\newcommand{\F}{\mathbb{F}}
\newcommand{\Fp}{\mathbb{F}_p}
\newcommand{\A}{\mathbb{A}}
\newcommand{\PP}{\mathbb{P}}
\newcommand{\Gal}{\mathrm{Gal}}
\newcommand{\End}{\mathrm{End}}
\newcommand{\Hom}{\mathrm{Hom}}
\newcommand{\Aut}{\mathrm{Aut}}
\newcommand{\id}{\mathrm{id}}
\newcommand{\NS}{\mathrm{NS}}
\newcommand{\Pic}{\mathrm{Pic}}
\newcommand{\CH}{\mathrm{CH}}
\newcommand{\cl}{\mathrm{cl}}
\newcommand{\Num}{\mathrm{Num}}
\DeclareMathOperator{\rank}{rank}
\DeclareMathOperator{\tr}{Spur}
\DeclareMathOperator{\codim}{kodim}
\DeclareMathOperator{\Spec}{Spec}

% -------------------------------------------------------
\begin{document}
% -------------------------------------------------------

\title{Grothendieck\textquoterightsche Standardvermutungen \"{u}ber algebraische Zyklen:\\
       Weil-Kohomologie, Motive und der Lefschetz-Operator}

\author{Michael Fuhrmann}
\address{Specialist Mathematics Project, Build 170}
\email{specialist-maths@localhost}
\date{\today}

\begin{abstract}
Grothendiecks vier Standardvermutungen --- Vermutung~A (Lefschetz),
Vermutung~B (Hodge-Standard), Vermutung~C (K\"unneth) und Vermutung~D
(Numerisch gleich homologisch) --- bilden die axiomatischen Grundpfeiler
der Theorie reiner Motive.  Ende der 1960er Jahre als algebraische Analoga
klassischer Sätze der Hodge-Theorie über $\mathbb{C}$ formuliert, sind diese
Vermutungen für Varietäten über beliebigen Körpern, insbesondere in positiver
Charakteristik, nach wie vor nahezu vollständig unbewiesen.  Wir geben
vollständige, strenge Formulierungen aller vier Vermutungen im Rahmen
der Weil-Kohomologietheorien, arbeiten ihre gegenseitigen Implikationen heraus,
stellen die Bezüge zu den Weil-Vermutungen (bewiesen von Deligne 1974),
zur Hodge-Vermutung und zur Tate-Vermutung dar und fassen die bekannten
Teilergebnisse zusammen: Kleimans Beweis, dass~D aus~A folgt, Jannsens
Halbeinfachheitssatz für numerische Motive, Andrés Theorie der motivierten Zyklen
sowie den aktuellen Stand der Forschung.
\end{abstract}

\maketitle

\tableofcontents

% -------------------------------------------------------
\section{Einleitung}
% -------------------------------------------------------

Das zentrale Ziel der algebraischen Geometrie ist das Studium algebraischer
Varietäten durch kohomologische Invarianten.  Über den komplexen Zahlen
gipfelt dieses Programm in der klassischen Hodge-Theorie: der harte
Lefschetz-Satz, die Hodge-Zerlegung und die positive Definitheit der
Weil-Form liefern eine reiche Struktur von $H^*(X,\Q)$ für eine glatte
projektive komplexe Mannigfaltigkeit~$X$.

Grothendiecks visionäre Einsicht, angekündigt auf dem Internationalen
Mathematikerkongress 1968 in Moskau und ausgearbeitet in seinem Briefwechsel
mit Serre und anderen Manuskripten, bestand darin, dass die tiefsten
Eigenschaften algebraischer Zyklen --- nicht nur über $\C$, sondern über
\emph{jedem} Körper $k$ --- durch eine kleine Anzahl präziser Vermutungen
erfasst werden sollten, die die klassische Hodge-Theorie widerspiegeln.
Die vier \emph{Standardvermutungen} tragen die Bezeichnungen A, B, C, D
(in Grothendiecks Originalnotation auch~I,~II für A und B).

\subsection*{Gliederung}

Abschnitt~\ref{sec:weil} erläutert die Axiome einer Weil-Kohomologietheorie.
Abschnitt~\ref{sec:zyklen} behandelt algebraische Zyklen und ihre Äquivalenzrelationen.
Abschnitt~\ref{sec:lefschetz} führt den Lefschetz-Operator und den klassischen
harten Lefschetz-Satz über $\C$ ein.
Die Abschnitte~\ref{sec:vermutungA}--\ref{sec:vermutungD} geben vollständige,
rigorose Formulierungen der Vermutungen~A--D.
Abschnitt~\ref{sec:implikationen} arbeitet die logischen Abhängigkeiten
zwischen den vier Vermutungen heraus.
Abschnitt~\ref{sec:weilverw} stellt den Bezug zu den Weil-Vermutungen und
Delignes Satz dar.
Abschnitt~\ref{sec:hodge} diskutiert den Zusammenhang mit der Hodge- und
Tate-Vermutung.
Abschnitt~\ref{sec:motive} erklärt die motivischen Konsequenzen.
Abschnitt~\ref{sec:andre} schildert Andrés motivierte Zyklen.
Abschnitt~\ref{sec:jannsen} präsentiert Jannsens Halbeinfachheitssatz.
Abschnitt~\ref{sec:bekannt} sammelt alle bekannten Spezialfälle.
Abschnitt~\ref{sec:hindernisse} erläutert die wesentlichen Hindernisse
für einen allgemeinen Beweis.
Abschnitt~\ref{sec:ausblick} gibt einen Ausblick auf offene Fragen.

% -------------------------------------------------------
\section{Weil-Kohomologietheorien}\label{sec:weil}
% -------------------------------------------------------

Im Folgenden ist $k$ ein Körper (oft ein endlicher Körper $\Fp$, ein
Zahlkörper oder $\C$), und alle Varietäten sind glatt, projektiv und
geometrisch zusammenhängend über $k$, sofern nicht anders angegeben.

\begin{definition_de}[Weil-Kohomologietheorie]
Eine \emph{Weil-Kohomologietheorie mit Koeffizientenkörper $F$} (von
Charakteristik null) ist ein kontravarianter Funktor $H^*$ von der
Kategorie glatter projektiver $k$-Varietäten in graduierte $F$-Algebren,
zusammen mit einer Zyklenklassenabbildung, der folgende Axiome erfüllt:
\begin{enumerate}[label=(\roman*)]
  \item \textbf{Endliche Dimensionalität.}  $H^i(X) = 0$ für $i < 0$ oder
        $i > 2\dim X$, und jedes $H^i(X)$ ist ein endlich-dimensionaler
        $F$-Vektorraum.
  \item \textbf{Poincaré-Dualität.}  Es gibt einen kanonischen Isomorphismus
        $H^{2d}(X)(d) \cong F$ (die \emph{Fundamentalklasse}), und das durch
        das Cup-Produkt induzierte Paar
        $H^i(X) \otimes H^{2d-i}(X) \to F$ ist eine perfekte Paarung.
        Dabei ist $d = \dim X$.
  \item \textbf{Künneth-Formel.}  Es gibt einen natürlichen Isomorphismus
        $H^*(X \times Y) \cong H^*(X) \otimes_F H^*(Y)$.
  \item \textbf{Zyklenklassenabbildung.}  Für jedes $i$ gibt es einen
        $\Q$-linearen Homomorphismus
        $\cl: \CH^i(X) \to H^{2i}(X)(i)$
        von der Chow-Gruppe der Kodimension-$i$-Zyklen, verträglich mit
        eigentlichem Pushforward, flachem Pullback und Cup-Produkt.
  \item \textbf{Schwacher Lefschetz / Hyperebenenabschnitt.}  Ist
        $Y \hookrightarrow X$ ein glatter Hyperebenenabschnitt, so ist
        $H^i(X) \to H^i(Y)$ ein Isomorphismus für $i < \dim Y$ und
        injektiv für $i = \dim Y$.
\end{enumerate}
\end{definition_de}

\begin{example}[Standardbeispiele von Weil-Kohomologietheorien]
\begin{enumerate}[label=(\alph*)]
  \item \textbf{Singuläre Kohomologie.}  Für $k \hookrightarrow \C$:
        $H^i(X) = H^i(X(\C),\Q)$ mit $F = \Q$.
  \item \textbf{De-Rham-Kohomologie.}  Für $k$ mit Charakteristik null:
        $H^i_{\mathrm{dR}}(X/k)$ mit $F = k$.
  \item \textbf{$\ell$-adische Kohomologie.}  Für $\ell \neq \ch(k)$:
        $H^i_{\mathrm{\acute{e}t}}(X_{\bar k}, \Q_\ell)$ mit $F = \Q_\ell$.
  \item \textbf{Kristalline Kohomologie.}  Für $k = \Fp$:
        $H^i_{\mathrm{kris}}(X/W(k))[\frac{1}{p}]$ mit $F = W(k)[\frac{1}{p}]$,
        nach Berthelot.
\end{enumerate}
\end{example}

\begin{bemerkung}
Alle bekannten Weil-Kohomologietheorien erfüllen die obigen Axiome.
Ein wesentliches Merkmal ist, dass sie für Varietäten, auf denen mehrere
Theorien gleichzeitig definiert sind, dieselben Betti-Zahlen $b_i(X) = \dim_F H^i(X)$
liefern.  Ob die vollen Kohomologiealgebren isomorph sind, erfordert jedoch die
Standardvermutungen.
\end{bemerkung}

% -------------------------------------------------------
\section{Algebraische Zyklen und Äquivalenzrelationen}\label{sec:zyklen}
% -------------------------------------------------------

Sei $X$ eine glatte projektive Varietät der Dimension $d$ über $k$.

\begin{definition_de}[Algebraischer Zyklus]
Ein \emph{algebraischer Zyklus} der Kodimension $p$ auf $X$ ist eine
endliche formale $\Z$-Linearkombination $Z = \sum n_i [V_i]$ irreduzibler
abgeschlossener Untervarietäten $V_i \subset X$ der Kodimension $p$.
Die Gruppe solcher Zyklen wird mit $\mathcal{Z}^p(X)$ bezeichnet.
\end{definition_de}

\begin{definition_de}[Äquivalenzrelationen auf Zyklen]
\begin{enumerate}[label=(\roman*)]
  \item \textbf{Rationale Äquivalenz}: $Z \sim_{\mathrm{rat}} 0$, wenn $Z$
        der Divisor einer rationalen Funktion auf einer Untervarietät ist.
        Der Quotient $\CH^p(X) = \mathcal{Z}^p(X)/{\sim_{\mathrm{rat}}}$
        ist die \emph{Chow-Gruppe}.
  \item \textbf{Algebraische Äquivalenz}: $Z \sim_{\mathrm{alg}} 0$, wenn
        $Z$ algebraisch äquivalent zu null ist (parametrisiert durch eine
        zusammenhängende Kurve).
  \item \textbf{Homologische Äquivalenz}: $Z \sim_{\mathrm{hom}} 0$, wenn
        $\cl(Z) = 0$ in $H^{2p}(X)(p)$ für eine (nach Vermutung~D: alle)
        Weil-Kohomologietheorien.
  \item \textbf{Numerische Äquivalenz}: $Z \sim_{\mathrm{num}} 0$, wenn
        $(Z \cdot W) = 0$ für alle Zyklen $W$ komplementärer Kodimension.
\end{enumerate}
Es gelten die Inklusionen
${\sim_{\mathrm{rat}}} \supseteq {\sim_{\mathrm{alg}}} \supseteq
 {\sim_{\mathrm{hom}}} \supseteq {\sim_{\mathrm{num}}}$.
\end{definition_de}

\begin{bemerkung}
Die letzte Inklusion folgt, weil die Schnittform durch die Zyklenklassenabbildung
faktorisiert.  Die Aussage ${\sim_{\mathrm{hom}}} = {\sim_{\mathrm{num}}}$
ist genau Vermutung~D.
\end{bemerkung}

% -------------------------------------------------------
\section{Der Lefschetz-Operator und der harte Lefschetz-Satz}\label{sec:lefschetz}
% -------------------------------------------------------

Sei $\xi \in \CH^1(X)$ die Klasse eines Hyperebenenabschnitts (d.h.\ eines
reichlich vorhandenen Geradenbündels $L$).  Seine Zyklenklasse ist
$L = \cl(\xi) \in H^2(X)(1)$.

\begin{definition_de}[Lefschetz-Operator]
Der \emph{Lefschetz-Operator} ist die Abbildung
\[
  \mathbf{L} = \cup L : H^i(X) \longrightarrow H^{i+2}(X)(1),
  \quad \alpha \mapsto L \cup \alpha.
\]
Seine iterierten Potenzen sind $\mathbf{L}^r : H^i(X) \to H^{i+2r}(X)(r)$.
\end{definition_de}

\begin{satz}[Harter Lefschetz-Satz über $\mathbb{C}$, Hodge 1950]
Sei $X$ eine glatte projektive komplexe Varietät der Dimension $d$, und
sei $L \in H^2(X,\Q)$ die Klasse eines Hyperebenenabschnitts.
Dann ist für alle $0 \leq i \leq d$:
\[
  \mathbf{L}^{d-i} : H^i(X,\Q) \xrightarrow{\;\sim\;} H^{2d-i}(X,\Q)
\]
ein Isomorphismus.
\end{satz}

Dieser Satz, klassisch über $\C$, ist über allgemeinen Körpern
\emph{unbekannt}.  Sein algebraisches Analogon ist der Inhalt von Vermutung~B.

\begin{definition_de}[Primitive Kohomologie und Lefschetz-Zerlegung]
Für $i \leq d$ ist die \emph{primitive Kohomologie}
$H^i_{\mathrm{prim}}(X) = \ker(\mathbf{L}^{d-i+1}: H^i(X) \to H^{2d-i+2}(X))$.
Die Lefschetz-Zerlegung (über $\C$) lautet:
\[
  H^i(X) = \bigoplus_{r \geq 0} \mathbf{L}^r H^{i-2r}_{\mathrm{prim}}(X).
\]
\end{definition_de}

% -------------------------------------------------------
\section{Vermutung~A: Die Lefschetz-Vermutung}\label{sec:vermutungA}
% -------------------------------------------------------

Vermutung~A behauptet, dass der Operator $\mathbf{L}^{d-i}$ ``algebraisch''
ist, d.h.\ durch einen algebraischen Zyklus (eine Korrespondenz) induziert wird.

\begin{vermutung}[Standardvermutung~A, Grothendieck 1969]\label{verm:A}
Sei $X$ eine glatte projektive Varietät der Dimension $d$ über einem Körper $k$,
und sei $H^*$ eine Weil-Kohomologietheorie mit Koeffizientenkörper $F$.
Für jedes $0 \leq i \leq d$ ist der Isomorphismus
\[
  \mathbf{L}^{d-i} : H^i(X) \xrightarrow{\;\sim\;} H^{2d-i}(X)(d-i)
\]
gegeben durch iteriertes Cup-Produkt mit $L = \cl(\xi)$, algebraisch,
d.h.\ er wird durch eine algebraische Korrespondenz $\Pi_i \in \CH^d(X \times X)$
induziert (modulo dem Twist).
Äquivalent: Die Projektoren $\pi_i$ in der Künneth-Zerlegung
$[\Delta_X] = \sum_{i=0}^{2d} \pi_i$ sind für alle $i$ algebraisch.
\end{vermutung}

\begin{bemerkung}
Vermutung~A hat zwei äquivalente Formulierungen:
\begin{enumerate}[label=(\alph*)]
  \item Der harte Lefschetz-Isomorphismus $\mathbf{L}^{d-i}$ ist algebraisch.
  \item Die Künneth-Komponenten $\pi_i$ der Diagonalklasse sind algebraisch.
\end{enumerate}
Die Äquivalenz $(a) \Leftrightarrow (b)$ ist eine formale Konsequenz der
Zerlegung der Identität; vgl.\ Kleiman~\cite{kleiman1968}.
\end{bemerkung}

\begin{satz}[Grothendieck]
Vermutung~A impliziert Vermutung~C.
\end{satz}

\begin{example}[Vermutung~A für Kurven]
Für eine glatte projektive Kurve $X$ ($d=1$) ist $\mathbf{L}^0 = \id$
trivial algebraisch.  Vermutung~A gilt also für Kurven.  Für Flächen ($d=2$)
reduziert sich Vermutung~A auf die Algebraizität der Lefschetz-Involution
$\Lambda: H^2(X) \to H^2(X)$.
\end{example}

% -------------------------------------------------------
\section{Vermutung~B: Die Hodge-Standardvermutung}\label{sec:vermutungB}
% -------------------------------------------------------

Vermutung~B ist in vielerlei Hinsicht die fundamentalste der vier, da sie
Positivitätseigenschaften der Schnittform impliziert.

\begin{definition_de}[Hodge-Riemann-Bilinearform]
Für primitive Klassen $\alpha, \beta \in H^i_{\mathrm{prim}}(X)$ sei
\[
  Q(\alpha, \beta) = (-1)^{i(i-1)/2} \int_X \alpha \cup \beta \cup L^{d-i}.
\]
Dies ist die \emph{Hodge-Riemann-Bilinearform}.
\end{definition_de}

\begin{vermutung}[Standardvermutung~B, Grothendieck 1969]\label{verm:B}
Sei $X$, $H^*$, $d$, $\xi$ wie oben.  Dann gilt:
\begin{enumerate}[label=(\roman*)]
  \item Die Lefschetz-Involution $\Lambda: H^i(X) \to H^{i-2}(X)$
        ist algebraisch, d.h.\ wird durch eine algebraische Korrespondenz
        induziert.
  \item Die Hodge-Riemann-Bilinearform $Q$ ist \emph{positiv definit} auf
        jedem Gradstück der primitiven Kohomologie.
\end{enumerate}
Äquivalent: Die $\mathfrak{sl}_2$-Wirkung, erzeugt von $\mathbf{L}$ und
$\Lambda$ auf $H^*(X)$, ist algebraisch.
\end{vermutung}

\begin{bemerkung}
Über $\C$ ist Vermutung~B eine Konsequenz der klassischen Hodge-Theorie
(der Hodge-Riemann-Bilinearrelationen).  Über beliebigen Körpern ist
Vermutung~B für Varietäten der Dimension $\geq 2$ (außer in Spezialfällen)
vollständig offen.
\end{bemerkung}

\begin{proposition}
Standardvermutung~B impliziert Standardvermutung~A.
\end{proposition}

\begin{satz}[Grothendieck]\label{satz:B_impliziert_D}
Standardvermutung~B (zusammen mit Vermutung~A) impliziert Standardvermutung~D.
\end{satz}

% -------------------------------------------------------
\section{Vermutung~C: Die Künneth-Vermutung}\label{sec:vermutungC}
% -------------------------------------------------------

\begin{definition_de}[Künneth-Zerlegung der Diagonale]
Aus der Künneth-Formel für $H^*$ ergibt sich eine Zerlegung
\[
  [\Delta_X] = \sum_{i=0}^{2d} \pi_i \in \bigoplus_{i=0}^{2d}
  H^{2d-i}(X) \otimes H^i(X) = H^{2d}(X \times X),
\]
wobei $[\Delta_X]$ die Klasse der Diagonale $\Delta_X \subset X \times X$ ist
und $\pi_i$ die Komponente in $H^{2d-i}(X) \otimes H^i(X)$ bezeichnet.
\end{definition_de}

\begin{vermutung}[Standardvermutung~C, Grothendieck 1969]\label{verm:C}
Für jedes $0 \leq i \leq 2d$ ist die Künneth-Komponente $\pi_i$ der
Diagonalklasse $[\Delta_X] \in H^{2d}(X \times X)(d)$ algebraisch:
\[
  \pi_i \in \mathrm{Bild}(\cl: \CH^d(X \times X) \otimes_\Z \Q \to
  H^{2d}(X \times X)).
\]
\end{vermutung}

\begin{bemerkung}
Vermutung~C ist äquivalent zur Existenz \emph{algebraischer} Projektoren
$\pi_i \in \CH^d(X \times X)_\Q$ mit $H^*(X) = \bigoplus_i \pi_i H^*(X)$.
Dies ist eine notwendige Bedingung dafür, dass $X$ eine
\emph{Chow-Künneth-Zerlegung} im Sinne von Murre besitzt.
\end{bemerkung}

\begin{proposition}[Eigenschaften von Vermutung~C]
\begin{enumerate}[label=(\roman*)]
  \item Vermutung~C wird von Vermutung~A impliziert.
  \item Über $\C$ sind die Künneth-Projektoren zwar durch die Hodge-Zerlegung
        gegeben, aber diese sind im Allgemeinen \emph{nicht} algebraisch ---
        Vermutung~C fragt genau danach.
  \item Für Varietäten mit einer zellulären Zerlegung (z.B.\ Fahnenvarietäten,
        Grassmannsche Varietäten) gilt Vermutung~C durch explizite Konstruktion.
\end{enumerate}
\end{proposition}

% -------------------------------------------------------
\section{Vermutung~D: Numerisch gleich homologisch}\label{sec:vermutungD}
% -------------------------------------------------------

Vermutung~D ist vielleicht die ``erdnäheste'' der vier, da sie eine rein
zykeltheoretische Aussage betrifft.

\begin{vermutung}[Standardvermutung~D, Grothendieck 1969]\label{verm:D}
Für jede glatte projektive Varietät $X$ über $k$ und für jede Primzahl
$\ell \neq \ch(k)$ (mit $\ell$-adischer Kohomologie) gilt:
\[
  Z \sim_{\mathrm{hom}} 0 \iff Z \sim_{\mathrm{num}} 0,
  \quad \forall Z \in \mathcal{Z}^p(X).
\]
Äquivalent: Die Zyklenklassenabbildung induziert eine Injektion
\[
  \cl: \CH^p(X)_\Q / {\sim_{\mathrm{num}}} \hookrightarrow H^{2p}(X)(p).
\]
\end{vermutung}

\begin{bemerkung}
Die Inklusion ${\sim_{\mathrm{hom}}} \supseteq {\sim_{\mathrm{num}}}$
ist elementar (Poincaré-Dualität).  Die nicht-triviale Behauptung ist die
Umkehrung: Wenn ein Zyklus jede Schnittform mit allen Zyklen komplementärer
Dimension null ergibt, dann ist seine Kohomologieklasse null.
\end{bemerkung}

\begin{satz}[Kleiman 1968]\label{satz:D_impliziert_A}
Standardvermutung~D impliziert Standardvermutung~A.
\end{satz}

\begin{proof}[Beweisidee]
Kleiman zeigte, dass die Injektivität der Zyklenklassenabbildung auf
numerischen Äquivalenzklassen erlaubt, die Künneth-Zerlegung der Diagonale
algebraisch durchzuführen.  Der Schlüssel liegt darin, dass der Künneth-Projektor
$\pi_i$ zusammen mit $\pi_j$ für $j \neq i$ numerisch trivial ist, und
Vermutung~D erzwingt dann homologische Trivialität auf dem Komplement,
was die Algebraizität liefert.  Vollständige Details in \cite{kleiman1968}.
\end{proof}

% -------------------------------------------------------
\section{Implikationen zwischen den Vermutungen}\label{sec:implikationen}
% -------------------------------------------------------

\begin{satz}[Implikationsdiagramm]\label{satz:implikationen}
Folgende Implikationen gelten unbedingt:
\[
  \mathbf{B} \Rightarrow \mathbf{A} \Rightarrow \mathbf{C},
  \qquad
  \mathbf{D} \Rightarrow \mathbf{A},
  \qquad
  \mathbf{A} + \mathbf{B} \Rightarrow \mathbf{D}.
\]
Insbesondere impliziert $\mathbf{B}$ alle vier Vermutungen gleichzeitig.
\end{satz}

\begin{proof}[Zusammenfassung der Beweise]
\begin{enumerate}[label=(\roman*)]
  \item $\mathbf{B} \Rightarrow \mathbf{A}$: Algebraizität von $\Lambda$ plus
        die $\mathfrak{sl}_2$-Relationen erzwingen die Algebraizität von
        $\mathbf{L}^{d-i}$.
  \item $\mathbf{A} \Rightarrow \mathbf{C}$: Die Künneth-Komponenten $\pi_i$
        werden durch die algebraischen Korrespondenzen für $\mathbf{L}^{d-i}$
        ausgedrückt.
  \item $\mathbf{D} \Rightarrow \mathbf{A}$: Kleimans Satz
        (Satz~\ref{satz:D_impliziert_A}).
  \item $\mathbf{A} + \mathbf{B} \Rightarrow \mathbf{D}$: Positive Definitheit
        (aus $\mathbf{B}$) plus algebraische Projektoren (aus $\mathbf{A}$)
        erzwingen, dass jeder homologisch triviale Zyklus numerisch trivial ist.
\end{enumerate}
\end{proof}

\begin{korollar}
Es ist nicht bekannt, ob $\mathbf{C}$ oder $\mathbf{D}$ allein die
Vermutung~$\mathbf{B}$ impliziert.  Die Implikation $\mathbf{D} \Rightarrow \mathbf{B}$
ist offen.
\end{korollar}

% -------------------------------------------------------
\section{Bezug zu den Weil-Vermutungen}\label{sec:weilverw}
% -------------------------------------------------------

Die Weil-Vermutungen (bewiesen von Deligne 1974 \cite{deligne1974}) besagen,
dass für eine glatte projektive Varietät $X$ über $\Fp$ die Zetafunktion
\[
  Z(X, T) = \exp\!\left(\sum_{n=1}^\infty |X(\mathbb{F}_{p^n})| \frac{T^n}{n}\right)
\]
folgendes erfüllt:
\begin{enumerate}[label=(\roman*)]
  \item \textbf{Rationalität}: $Z(X,T) \in \Q(T)$.
  \item \textbf{Funktionalgleichung}: $Z(X, q^{-d}T^{-1}) = \pm q^{d\chi/2} T^\chi Z(X,T)$.
  \item \textbf{Riemann-Hypothese}: Die Eigenwerte des Frobenius auf $H^i$
        haben Betrag $p^{i/2}$.
  \item \textbf{Betti-Zahlen}: $Z(X,T) = \prod_{i=0}^{2d} P_i(T)^{(-1)^{i+1}}$
        mit $\deg P_i = b_i(X)$.
\end{enumerate}

\begin{satz}[Grothendiecks bedingter Beweis]
Unter Annahme der Standardvermutungen (insbesondere Vermutung~B) kann man
den Riemann-Hypothese-Teil der Weil-Vermutungen beweisen.
\end{satz}

\begin{bemerkung}
Deligne fand einen völlig anderen Beweis der Weil-Vermutungen, der die
Standardvermutungen \emph{nicht} verwendet.  Dies macht die Standardvermutungen
noch geheimnisvoller: Sie implizieren mehr über die Struktur von Zyklen
und Motiven, als für das bloße Zählen von Punkten nötig wäre.
\end{bemerkung}

% -------------------------------------------------------
\section{Bezug zur Hodge- und Tate-Vermutung}\label{sec:hodge}
% -------------------------------------------------------

\begin{vermutung}[Hodge-Vermutung, Hodge 1950]
Sei $X$ eine glatte projektive komplexe Varietät.  Dann ist jede rationale
Hodge-Klasse in $H^{2p}(X,\Q) \cap H^{p,p}(X)$ eine rationale Linearkombination
von Zyklenklassen $\cl(Z)$ für algebraische Zyklen $Z \in \mathcal{Z}^p(X)$.
\end{vermutung}

\begin{proposition}[Hodge impliziert A über $\C$]
Die Hodge-Vermutung (für $X$ und $X \times X$) impliziert Standardvermutung~A
über $\C$.
\end{proposition}

\begin{vermutung}[Tate-Vermutung, Tate 1963]
Sei $X$ eine glatte projektive Varietät über einem endlich erzeugten Körper $k$,
und $\ell$ eine Primzahl verschieden von $\ch(k)$.  Dann ist die Zyklenklassenabbildung
\[
  \cl: \CH^p(X)_{\Q_\ell} \to H^{2p}(X_{\bar k}, \Q_\ell(p))^{G_k}
\]
surjektiv (wobei $G_k = \Gal(\bar k/k)$).
\end{vermutung}

\begin{satz}[Tate impliziert D über endlichen Körpern]
Die Tate-Vermutung (Surjektivität) sowie ihre Injektivitätsvariante impliziert
Standardvermutung~D über endlichen Körpern.
\end{satz}

% -------------------------------------------------------
\section{Motivische Konsequenzen}\label{sec:motive}
% -------------------------------------------------------

\begin{definition_de}[Kategorie reiner Motive]
Sei $\sim$ eine geeignete Äquivalenzrelation auf algebraischen Zyklen.
Die \emph{Kategorie reiner Motive} $\mathcal{M}_\sim(k)$ hat:
\begin{itemize}
  \item \textbf{Objekte}: Paare $(X, p)$, wobei $X$ glatt projektiv und
        $p \in \CH^d(X \times X)_\Q$ ein idempotenter Korrespondenz-Projektor ist.
  \item \textbf{Morphismen}: $\Hom((X,p),(Y,q)) = q \circ \CH^d(X \times Y)_\Q \circ p$.
  \item \textbf{Tensorprodukt}: $(X,p) \otimes (Y,q) = (X \times Y, p \times q)$.
\end{itemize}
\end{definition_de}

\begin{satz}[Jannsen 1992]\label{satz:jannsen}
Die Kategorie $\mathcal{M}_{\mathrm{num}}(k)$ reiner Motive modulo numerischer
Äquivalenz ist eine \emph{halbeinfache abelsche Kategorie}.
\end{satz}

Dieser Satz, von Jannsen \emph{ohne} Annahme der Standardvermutungen bewiesen,
ist eines der großen unbedingten Ergebnisse der Motivtheorie.

\begin{satz}[Bedingter Struktursatz für Motive]
Unter Annahme der Standardvermutungen zerfällt der Funktor
$h: X \mapsto (X, [\Delta_X])$ als
$h(X) = \bigoplus_{i=0}^{2d} h^i(X)$,
wobei jedes $h^i(X)$ ein reines Motiv vom Gewicht $i$ ist.
Ferner ist $\mathcal{M}_{\mathrm{num}}(k)$ eine Tannaka-Kategorie.
\end{satz}

% -------------------------------------------------------
\section{Andrés motivierte Zyklen}\label{sec:andre}
% -------------------------------------------------------

In Ermangelung eines Beweises der Standardvermutungen führte
Andr\'e~\cite{andre1996} eine größere Klasse ``motivierter Zyklen'' ein.

\begin{definition_de}[Motivierte Zyklen, Andr\'{e} 1996]
Eine Kohomologieklasse $\alpha \in H^{2p}(X)(p)$ heißt \emph{motivierter
Zyklus}, wenn sie zur kleinsten $\Q$-Unteralgebra von $\bigoplus_X H^*(X)$
gehört, die von folgendem erzeugt wird:
\begin{enumerate}[label=(\roman*)]
  \item Algebraischen Zyklenklassen $\cl(\mathcal{Z}^p(X))$.
  \item Klassen der Form $*_L(\beta)$ für motivierte $\beta$, wobei $*_L$
        der Hodge-$*$-Operator zu einer Polarisation $L$ ist.
\end{enumerate}
\end{definition_de}

\begin{satz}[Andr\'{e} 1996]\label{satz:andre}
Die Kategorie $\mathcal{M}_{\mathrm{mot}}(k)$ der aus motivierten Zyklen
gebauten Motive ist:
\begin{enumerate}[label=(\roman*)]
  \item Eine halbeinfache $\Q$-lineare abelsche Tensorkategorie.
  \item Erfüllt alle vier Standardvermutungen (per Konstruktion).
  \item Enthält die abelschen Motive als Unterkategorie.
\end{enumerate}
\end{satz}

\begin{bemerkung}
Andrés Satz beweist die Standardvermutungen \emph{nicht}: Er zeigt nur,
dass eine Erweiterung der algebraischen Zyklen um den $*_L$-Operator die
gewünschten Eigenschaften hat.  Die Standardvermutungen fragen, ob $*_L$
bereits algebraisch ist.
\end{bemerkung}

% -------------------------------------------------------
\section{Jannsens Halbeinfachheitssatz}\label{sec:jannsen}
% -------------------------------------------------------

\begin{satz}[Jannsen 1992, vollständige Aussage]\label{satz:jannsen_voll}
Sei $k$ ein beliebiger Körper.  Die Kategorie $\mathcal{M}_{\mathrm{num}}(k)$
reiner Motive modulo numerischer Äquivalenz (über $k$, mit $\Q$-Koeffizienten)
ist eine halbeinfache abelsche rigide Tensorkategorie.
\end{satz}

\begin{proof}[Beweisidee]
Die Hauptschritte sind:
\begin{enumerate}[label=(\arabic*)]
  \item Zeige, dass $\End_{\mathcal{M}_{\mathrm{num}}}(h(X))$ für alle $X$
        eine endlich-dimensionale $\Q$-Algebra ist.
  \item Zeige, dass diese Algebra keine Nilpotente hat (unter Verwendung der
        Tatsache, dass die numerische Äquivalenz die ``gröbste'' geeignete
        Äquivalenzrelation ist).
  \item Schließe auf Halbeinfachheit durch ein allgemeines kategorientheoretisches
        Argument.
\end{enumerate}
\end{proof}

% -------------------------------------------------------
\section{Bekannte Ergebnisse und Spezialfälle}\label{sec:bekannt}
% -------------------------------------------------------

\begin{satz}[Bekannte Fälle von Vermutung~A]
Standardvermutung~A ist bekannt für:
\begin{enumerate}[label=(\roman*)]
  \item Glatte projektive Kurven ($d=1$): trivial.
  \item Glatte projektive Flächen ($d=2$): folgt aus dem Satz von Lefschetz
        über $(1,1)$-Klassen.
  \item Abelsche Varietäten: Lieberman~\cite{lieberman1968} bewies Vermutung~A
        mithilfe der Rosati-Involution.
  \item Varietäten mit zellulärer Zerlegung (z.B.\ Fahnenvarietäten,
        Grassmannsche): explizite Konstruktion der Künneth-Projektoren.
\end{enumerate}
\end{satz}

\begin{satz}[Bekannte Fälle von Vermutung~B]
Standardvermutung~B ist bekannt für:
\begin{enumerate}[label=(\roman*)]
  \item Kurven und Flächen über $\C$: klassische Hodge-Theorie.
  \item Abelsche Varietäten über beliebigen Körpern: Rosati-Involution
        (vgl.\ Mumford~\cite{mumford1970}).
  \item $K3$-Flächen: verifiziert von Kleiman u.a.
\end{enumerate}
Vermutung~B ist \emph{offen} für alle Varietäten der Dimension $\geq 3$
über Körpern positiver Charakteristik.
\end{satz}

\begin{satz}[Bekannte Fälle von Vermutung~D]
Standardvermutung~D ist bekannt für:
\begin{enumerate}[label=(\roman*)]
  \item Divisoren auf beliebigen glatten projektiven Varietäten (N\'eron-Severi-Theorem).
  \item Nullzyklen auf Flächen: Satz von Bloch-Srinivas~\cite{bloch1983}.
  \item Abelsche Varietäten: aus der Arbeit von Lieberman u.a.
\end{enumerate}
\end{satz}

% -------------------------------------------------------
\section{Hindernisse für einen allgemeinen Beweis}\label{sec:hindernisse}
% -------------------------------------------------------

\begin{bemerkung}[Warum die Standardvermutungen schwer sind]
Die grundlegende Schwierigkeit ist die \emph{Transzendenz} des Hodge-$*$-Operators
$*_L$.  In der klassischen Hodge-Theorie über $\C$ wird $*_L$ durch die
Metrik der Varietät definiert, die keine algebraische Interpretation hat.
Die Standardvermutungen fordern einen rein algebraischen Ersatz.

Insbesondere:
\begin{enumerate}[label=(\roman*)]
  \item Es gibt keine bekannte algebraische Konstruktion von $\Lambda$, die
        in positiver Charakteristik funktioniert.
  \item Die Ansätze über Witt-Kohomologie und kristalline Kohomologie haben
        bisher keinen Beweis von Vermutung~B geliefert.
  \item Das Langlands-Programm gibt keinen direkten Zugang zu Zyklen.
  \item Über endlichen Körpern sind Motive besser verstanden (via Frobenius),
        aber auch dort bleibt Vermutung~B für $d \geq 3$ offen.
\end{enumerate}
\end{bemerkung}

% -------------------------------------------------------
\section{Ausblick und offene Fragen}\label{sec:ausblick}
% -------------------------------------------------------

\begin{vermutung}[Murres Vermutungen, 1993]
Für jede glatte projektive Varietät $X$ der Dimension $d$ existieren
Idempotente $\Pi_0, \Pi_1, \ldots, \Pi_{2d} \in \CH^d(X \times X)_\Q$
mit $\sum_i \Pi_i = [\Delta_X]$, $\Pi_i \circ \Pi_j = \delta_{ij} \Pi_i$,
und $(\Pi_i)_*$ wirkt auf $H^j(X)$ als $\delta_{ij}\cdot\id$.
\end{vermutung}

\begin{bemerkung}[Das motivische Kohomologieprogramm]
Voevodskys Theorie der motivischen Kohomologie~\cite{voevodsky2000} bietet
einen Rahmen zum Verständnis algebraischer Zyklen, der über die Standardvermutungen
hinausgeht.  Insbesondere liefern die Milnor-Vermutung (bewiesen von Voevodsky)
und die Bloch-Kato-Vermutung (bewiesen von Voevodsky-Rost) tiefe Ergebnisse
über die Struktur der Milnor-$K$-Theorie, implizieren aber nicht direkt die
Standardvermutungen.
\end{bemerkung}

\begin{bemerkung}[Potenzielle Ansätze via derivierte Kategorien]
Bondal-Orlows Rekonstruktionssatz~\cite{bondal2001} legt nahe, dass die
derivierte Kategorie $D^b(\mathrm{Koh}(X))$ alle motivischen Informationen
enthält.  Falls man derivierte Äquivalenzen mit motivierten Zyklen in Bezug
setzen könnte, ergäbe sich ein neuer Ansatz zu den Standardvermutungen.
\end{bemerkung}

\begin{thebibliography}{99}

\bibitem{grothendieck1968}
A.\ Grothendieck,
\textit{Standard conjectures on algebraic cycles},
in: Algebraic Geometry (Internat.\ Colloq., Tata Inst.\ Fund.\ Res., Bombay, 1968),
Oxford Univ.\ Press, London, 1969, S.\ 193--199.

\bibitem{kleiman1968}
S.\ L.\ Kleiman,
\textit{Algebraic cycles and the Weil conjectures},
in: Dix expos\'es sur la cohomologie des sch\'emas, Masson \& Cie, Paris;
North-Holland, Amsterdam, 1968, S.\ 359--386.

\bibitem{deligne1974}
P.\ Deligne,
\textit{La conjecture de Weil, I},
Publ.\ Math.\ IHES \textbf{43} (1974), 273--307.

\bibitem{deligne1980}
P.\ Deligne,
\textit{La conjecture de Weil, II},
Publ.\ Math.\ IHES \textbf{52} (1980), 137--252.

\bibitem{jannsen1992}
U.\ Jannsen,
\textit{Motives, numerical equivalence, and semi-simplicity},
Invent.\ Math.\ \textbf{107} (1992), 447--452.

\bibitem{andre1996}
Y.\ Andr\'e,
\textit{Pour une th\'eorie inconditionnelle des motifs},
Publ.\ Math.\ IHES \textbf{83} (1996), 5--49.

\bibitem{tate1994}
J.\ Tate,
\textit{Conjectures on algebraic cycles in $\ell$-adic cohomology},
in: Motives (Seattle, 1991), Proc.\ Sympos.\ Pure Math.\ \textbf{55}, Teil 1,
Amer.\ Math.\ Soc., Providence, RI, 1994, S.\ 71--83.

\bibitem{milne2002}
J.\ S.\ Milne,
\textit{Lefschetz motives and the Tate conjecture},
Compositio Math.\ \textbf{117} (1999), 45--76.

\bibitem{lieberman1968}
D.\ Lieberman,
\textit{Numerical and homological equivalence of algebraic cycles on Hodge manifolds},
Amer.\ J.\ Math.\ \textbf{90} (1968), 366--374.

\bibitem{mumford1970}
D.\ Mumford,
\textit{Abelian Varieties},
Tata Institute of Fundamental Research Studies in Mathematics,
Oxford Univ.\ Press, 1970.

\bibitem{bloch1983}
S.\ Bloch, V.\ Srinivas,
\textit{Remarks on correspondences and algebraic cycles},
Amer.\ J.\ Math.\ \textbf{105} (1983), 1235--1253.

\bibitem{lefschetz1924}
S.\ Lefschetz,
\textit{L'Analysis Situs et la G\'eom\'etrie Alg\'ebrique},
Gauthier-Villars, Paris, 1924.

\bibitem{voevodsky2000}
V.\ Voevodsky,
\textit{Triangulated categories of motives over a field},
in: Cycles, Transfers, and Motivic Homology Theories,
Ann.\ of Math.\ Stud.\ \textbf{143}, Princeton Univ.\ Press, 2000, S.\ 188--238.

\bibitem{bondal2001}
A.\ Bondal, D.\ Orlov,
\textit{Reconstruction of a variety from the derived category and groups of autoequivalences},
Compositio Math.\ \textbf{125} (2001), 327--344.

\bibitem{beilinson1985}
A.\ A.\ Beilinson,
\textit{Higher regulators and values of $L$-functions},
J.\ Soviet Math.\ \textbf{30} (1985), 2036--2070.

\bibitem{kleiman1994}
S.\ L.\ Kleiman,
\textit{The standard conjectures},
in: Motives (Seattle, 1991), Proc.\ Sympos.\ Pure Math.\ \textbf{55}, Teil 1,
Amer.\ Math.\ Soc., Providence, RI, 1994, S.\ 3--20.

\bibitem{murre1993}
J.\ P.\ Murre,
\textit{On a conjectural filtration on the Chow groups of an algebraic variety, I, II},
Indag.\ Math.\ (N.S.) \textbf{4} (1993), 177--201.

\bibitem{scholl1994}
A.\ J.\ Scholl,
\textit{Classical motives},
in: Motives (Seattle, 1991), Proc.\ Sympos.\ Pure Math.\ \textbf{55}, Teil 1,
Amer.\ Math.\ Soc., Providence, RI, 1994, S.\ 163--187.

\end{thebibliography}

\end{document}
